\section{The tool}
\label{sec:tool}

\paragraph{Files} (local mirror \fpath{~/Work/Hermes/sPHENIX/run3pp/trackqa/}, cluster
\fpath{/gpfs02/sphenix/user/dongxs/run3pp/trackqa/}):
\begin{itemize}\itemsep0pt
\item \fpath{src/TrackQATree/} -- the SubsysReco module (autotools package, libraries
  \fcmd{fun4all}, \fcmd{ffaobjects}, \fcmd{ffarawobjects}, \fcmd{trackbase_historic(_io)},
  \fcmd{globalvertex_io}); built by \fpath{build.sh} into the project install area shared with
  \fcmd{EventInfoTree}.
\item \fpath{macro/Fun4All_TrackQA.C(nEvents, input, outdir, trimDST, minPt)} -- DST input
  manager (single file or a text list of files of one run), the module, and optionally a
  \fcmd{Fun4AllDstOutputManager} for the trimmed DST. Prints \fcmd{All done} only when all
  outputs exist.
\item \fpath{macro/run_trackqa.sh <nEvents> <input> <outdir> [trimDST 0/1] [minPt]} -- job wrapper:
  clean environment, scratch directory, log with \fcmd{/usr/bin/time -v}, fails unless the log
  contains \fcmd{All done}. Same conventions as \fcmd{run_kfp.sh}, so the existing condor
  scripts can be adapted by changing the executable.
\item \fpath{tools/check_vertex_id.C} -- diagnostic used to define the vertex association
  (Sec.~\ref{sec:vtxassoc}).
\item \fpath{analysis_note/scripts/make_qa.C} -- figures, numbers and checks of this note (runs
  locally on the downloaded ntuple).
\end{itemize}

\choice{The job reads the DST nodes only: no CDB, geometry, field map or track refit. The DCA
calculation (\fcmd{TrackAnalysisUtils::get_dca}) needs only the stored track state and
covariance at the point of closest approach. This keeps the job at \cpuMsPerEv\,ms/event and
independent of calibration tags. Quantities that need geometry or clusters (dE/dx, residuals)
are therefore not in the ntuple.}

\choice{No event or track selection is applied when writing: every DST event has an
\fcmd{event} entry and every track of \fcmd{SvtxTrackMap} a \fcmd{track} entry
(\fcmd{minPt}=0 by default). Selections are cheap to apply offline, and the
counts can be checked against the DST (Sec.~\ref{sec:audit}). The \fcmd{minPt} argument exists
to shrink the track tree if needed; tracks below it are still counted in the event tree.}

\subsection{Track classes and crossing}
Following the KFParticle study, tracks are classified by their cluster content
(\fcmd{SvtxTrack_v5} cluster counters): \emph{INTT} ($\geq1$ INTT cluster, crossing measured by
INTT timing), \emph{MVTX-only} (MVTX but no INTT clusters) and \emph{TPC-only} (no silicon
cluster). The track beam crossing \fcmd{crossing} is relative to the GL1 BCO of the event; TPC-only
tracks always carry crossing 0 because no silicon timing is available (verified for every track,
check \fcmd{chkTpcOnlyXzero}).

\subsection{Vertex association for the DCA}
\label{sec:vtxassoc}
\fcmd{SvtxTrack::get_vertex_id()} was checked on 200 events with \fpath{tools/check_vertex_id.C}:
when set, it always points to an existing vertex of the track's own crossing; the vertex track
lists are consistent with the ids (every listed track points back to the vertex), but many tracks
carry the id of a vertex they were not used in. The ntuple therefore stores the class
\fcmd{vtxMatch}:
\begin{center}\small
\begin{tabular}{clr}
\toprule
\fcmd{vtxMatch} & definition & tracks (\%)\\
\midrule
1 & vertex = stored vertex id, track in that vertex's track list (used in the fit) & \nMatchOne{} (\fracMatchOne)\\
2 & vertex = stored vertex id, track not in the fit & \nMatchTwo{} (\fracMatchTwo)\\
3 & no vertex id; vertex of the same crossing closest in $z$ to the track PCA & \nMatchThree{} (\fracMatchThree)\\
0 & no vertex in the track's crossing; DCA and vertex branches $=-999$ & \nMatchZero{} (\fracMatchZero)\\
\bottomrule
\end{tabular}
\end{center}
In the test segment class 3 contains only TPC-only tracks and classes 1--2 only silicon tracks
(check \fcmd{chkTpcVtxId}).

\choice{DCA reference = the stored vertex id when present (the tracking's own assignment),
otherwise the closest vertex in $z$ of the same crossing. The class is stored, so analyses can
choose; e.g.\ \fcmd{vtxMatch==1} for primary-track QA.}

\caveat{For \fcmd{vtxMatch==1} the track was used in the vertex fit, so its DCA is biased towards
zero, strongly for vertices with few tracks (Sec.~\ref{sec:qa}). For TPC-only tracks (class 3) the
vertex is a crossing-0 vertex by construction, while most TPC-only tracks come from other
crossings (their $z$ is then shifted by the drift-time offset): their DCA is not a primary/secondary
discriminator.}
